\documentclass[acmsmall,review,prologue,dvipsnames]{acmart}
\settopmatter{printfolios=true,printccs=false,printacmref=true}

\usepackage{pig}
\ColourEpigram
\usepackage[T1]{fontenc}
\usepackage{hyperref}
\usepackage{amsmath}
\usepackage{stmaryrd}
\usepackage{upgreek}

%%%% ascii
\newcommand{\hash}{\char 35}

%%%% colour and font conventions
\newcommand{\remph}{\emph}  %% rhetorical emphasis, i.e., ranting
\newcommand{\demph}{\emph}  %% definitional emphasis, i.e., pointing out the defining occurrence of some locally new terminology
\newcommand{\D}[1]{\blue{\mathsf{#1}}}
\newcommand{\DU}[1]{\D{\mathbb{#1}}}
\newcommand{\V}[1]{\purple{\mathit{#1}}}
\newcommand{\C}[1]{\red{\mathsf{#1}}}
\newcommand{\F}[1]{\green{\mathbf{#1}}}

%%%% specific identifiers
\newcommand{\UF}{\DU{F}}
\newcommand{\ElF}[1]{\F{\llbracket \black{#1} \rrbracket}_{\UF}}

\newcommand{\UE}{\F{\mathbb{E}}}
\newcommand{\UD}{\DU{D}}
\newcommand{\SET}{\D{Set}}

\newcommand{\CTR}{\D{Container}}
\newcommand{\CX}[2]{\F{\llbracket \black{#1} \rrbracket^{\D{C}}}\:#2}
\newcommand{\pCX}[1]{(#1)}
\newcommand{\ctr}[2]{#1 \mathop{\red{\triangleleft}} #2}
\newcommand{\FTR}{\D{Fontainer}}
\newcommand{\FX}[2]{\F{\llbracket \black{#1} \rrbracket^{\D{F}}}\:#2}
\newcommand{\CSh}[1]{#1\:\F{.Sh}}
\newcommand{\CPo}[2]{#1\:\F{.Po}\:#2}

\newcommand{\FM}[2]{{#1}^{\D{\ast}}\:#2}
\newcommand{\va}[1]{\red{\hash}#1}
\newcommand{\nd}[1]{\red{\langle\black{#1}\rangle}}
\newcommand{\REC}[1]{\D{Rec}\:#1}
\newcommand{\rec}[1]{\F{rec}\:#1}

\newcommand{\fto}{\mathrel{\blue{\to}}}
\newcommand{\pti}{\mathrel{\blue{\times}}}
\newcommand{\peq}{\mathrel{\blue{=}}}
\newcommand{\blank}{\cdot}

\newcommand{\SG}[3]{(\V{#1}\hb #2)\pti #3}
\newcommand{\pr}[2]{#1\red{,}\, #2}

\newcommand{\PI}[3]{(\V{#1}\hb #2)\fto #3}
\newcommand{\LAM}[2]{\red{\uplambda}\V{#1}\mapsto #2}
\newcommand{\KON}[1]{\red{\upkappa}\:#1}

%%% Various combinators
\newcommand{\ENUM}[1]{\blue{\{}#1\blue{\}}}

\newcommand{\VERTPI}[2]{#1 \mathrel{{}_{\UF}\!\!\green{\raisebox{0.03in}[0in][0in]{\ensuremath{\to}}}} #2}
\newcommand{\vertflam}[1]{\F{\uplambda}_{\UF}\:#1}
\newcommand{\vertfapp}[2]{#1 \mathop{\F{\$}_{\UF}} #2}
\newcommand{\VERTPIF}[2]{#1 \mathrel{\green{\raisebox{0.03in}[0in][0in]{\ensuremath{\to}}}_{\UF}} #2}
\newcommand{\ROUGEPIF}[2]{#1 \mathrel{\red{\raisebox{0.03in}[0in][0in]{\ensuremath{\to}}}_{\UF}} #2}

\newcommand{\BLEUPI}[2]{#1 \mathrel{{}_{\UF}\!\!\blue{\raisebox{0.03in}[0in][0in]{\ensuremath{\to}}}} #2}
\newcommand{\VERTPIB}[2]{#1 \mathrel{{}_{\UF}\!\!\green{\raisebox{0.03in}[0in][0in]{\ensuremath{\leadsto}}}} #2}
\newcommand{\bleuflam}[1]{\F{\uplambda}_{\UF}\:#1}
\newcommand{\bleufapp}[2]{#1 \mathop{\F{\$}_{\UF}} #2}


\newcommand{\FSG}[2]{#1 \mathrel{\red{\times_{\UF}}} #2}
\newcommand{\impl}[1]{\ensuremath{\left\{ #1 \right\}}}
\newcommand{\HOLE}{\orange{?}}

%%%% punctuation and spacing
\newcommand{\hab}{:}
\newcommand{\hb}{\!\hab\!}  %% thinner for use in sigma and pi types
%\hg for "hypothesis gap" is in pig.sty

%%% spelling
\newcommand{\jgm}[1]{judgement#1}

%%%% meta
\newcommand{\TODO}[1]{\textbf{TODO: #1}}

\begin{document}
\title{Search Me!}

\author{Jacques Carette}
\affiliation{
\department{Department of Computing and Software}
\institution{McMaster University}
\city{Hamilton}
\country{Canada}
}
\author{Conor Titania Mc Bride}
\affiliation{
\department{Department of Computer and Information Sciences}
\institution{University of Strathclyde}
\city{Glasgow}
\country{Scotland}
}

%% Journal information
%% Supplied to authors by publisher for camera-ready submission;
%% use defaults for review submission.
\acmJournal{PACMPL}
\acmVolume{1}
\acmNumber{CONF} % CONF = POPL or ICFP or OOPSLA
\acmArticle{1}
\acmYear{2018}
\acmMonth{1}
\acmDOI{} % \acmDOI{10.1145/nnnnnnn.nnnnnnn}
\startPage{1}

\setcopyright{none}

%% Bibliography style
\bibliographystyle{ACM-Reference-Format}
%% Citation style
%% Note: author/year citations are required for papers published as an
%% issue of PACMPL.
\citestyle{acmauthoryear}   %% For author/year citations

%% 2012 ACM Computing Classification System (CSS) concepts
%% Generate at 'http://dl.acm.org/ccs/ccs.cfm'.
\begin{CCSXML}
<ccs2012>
<concept>
<concept_id>10011007.10011006.10011008</concept_id>
<concept_desc>Software and its engineering~General programming languages</concept_desc>
<concept_significance>500</concept_significance>
</concept>
<concept>
<concept_id>10003456.10003457.10003521.10003525</concept_id>
<concept_desc>Social and professional topics~History of programming languages</concept_desc>
<concept_significance>300</concept_significance>
</concept>
</ccs2012>
\end{CCSXML}

\ccsdesc[500]{Software and its engineering~General programming languages}
\ccsdesc[300]{Social and professional topics~History of programming languages}
%% End of generated code

%% Keywords
%% comma separated list
\keywords{keyword1, keyword2, keyword3}  %% \keywords are mandatory in final camera-ready submission

\begin{abstract}
Fill me in
\end{abstract}

\maketitle

\section{Introduction}\label{sec:intro}

Up to isomorphism, this paper is about numbers --- how many elements belong to types which are finite. However, we are anxious to retain the \remph{structure} of our finitary data and disinclined to \remph{count} them, so encoding finite types by their mere cardinality is asking for trouble. What matters more than the \remph{size} of a finite type is the exhaustive \remph{exploration} of its elements. We can achieve the latter with more sensitivity than the British method of forming an orderly queue.

Our prospectus is to construct two type theoretical universes, $\UF$ of \remph{finite data}, and $\UD$ of \remph{decisions}, with $\UD$ admitting universal and existential quantification over types in $\UF$. Decisions are thus undertaken by exhaustive exploration: what $\UF$-types say to $\UD$-types is `Search me!'.

The $\UF$ universe offers \remph{tagged} enumerations. Its dependent pairs are actual pairs, with projections that do not rely on the slightest arithmetic. Moreover, functions from $\UF$-types can be \remph{tabulated} as first-order data structures, ensuring that extensionally equal such functions are represented by intensionally equal tables.

Meanwhile, the $\UD$ universe is more logical in flavour, generated by quantification, but with notational conveniences for the usual nullary and binary connectives. Negation is \remph{admissible} by de Morgan duality, rather than by `implies false', meaning that the evidence for the negation of a universal statement is the constructive existence of a counterexample. Again, the searchability of the domains of quantification is what makes this possible. Decisions in $\UD$ are computable only because they are finitary, so it is crucial to give $\UD$ its meaning by interpretation into $\UF$. As a consequence $\UD$ comes with a notion of \remph{comprehension}.

What is it for? We often work with treelike (i.e., termlike) data structures with finitely many node shapes, and for each such, finitely many subtree positions. Certainly, Universal Algebra is concerned with such things, as is a great deal of perfectly ordinary functional programming. Finding these shape and position types in $\UF$ allows us to express properties of these trees in $\UD$, and thus decide them by construction.

\TODO{Say something about Mart\'\i{}n's much more subtle notions of searchability.}

\section{Free Monads of Containers}

By way of motivation, let us start where we aim to finish. We build termlike data as free monads of containers (in the sense of \citet{DBLP:journals/tcs/AbbottAG05}).
\[
\Rule{\V{C} \hab \CTR \hg  \V{X} \hab \SET}
     {\FM{\V{C}}{\V{X}} \hab \SET}
\qquad
\Rule{\V{x}\hab \V{X}}
     {\va{\V{x}}\hab \FM{\V{C}}{\V{X}}}
\qquad
\Rule{\V{cts} \hab \CX{\V{C}}{\pCX{\FM{\V{C}}{\V{X}}}}}
     {\nd{\V{cts}}\hab \FM{\V{C}}{\V{X}}}
\]
where a container is specified by a shape type and a shape-indexed family of positions, then interpreted as the dependent pair of a shape and a function from positions to payload:
\[
\Rule{\V{S}\hab\SET\hg \V{P}\hab\V{S}\fto\SET}
     {\ctr{\V{S}}{\V{P}}\hab\CTR}
\qquad
\Rule{\V{C}\hab\CTR\hg \V{X}\hab\SET}
     {\CX{\V{C}}{\V{X}}\hab\SET}
\quad
\CX{\ctr{\V{S}}{\V{P}}}{\V{X}} = \SG{s}{\V{S}}{(\V{P}\:\V{s}\fto\V{X})}
\]
Notationally, if it is inconvenient to pattern match a container, $\V{C}$, as $\ctr{\V{S}}{\V{P}}$, we may write its shape $\V{S}$ as $\CSh{\V{C}}$ and instantiate its position family $\V{P}\:\V{s}$ as $\CPo{\V{C}}{\V{s}}$.

To equip $\FM{\V{C}}{\blank}$ with the structure of a monad, we should note that $\va \blank$ is `return' and define `join' as follows:
\[ \Rule{\V{ts} \hab \FM{\V{C}}{(\FM{\V{C}}{\V{X}})}}
        {\F{join}\: \V{ts} \hab \FM{\V{C}}{\V{X}}}
\qquad\qquad
\begin{array}[c]{lcl}
\F{join}\: \va{\V{t}} & = & \V{t}\\
\F{join}\: \nd{\pr{\V{s}}{\V{k}}} & = & \nd{\pr{\V{s}}{\LAM{p}{\F{join}\: (\V{k}\:\V{p})}}} \\
\end{array}
\]
The recursion in this definition is structural, by inspection, because all of $\V{k}$'s output are smaller than $\nd{\pr{\V{s}}{\V{k}}}$. One of the monad laws \emph{is} the base case of the definition, but to prove the other two requires function extensionality.

Meanwhile, by way of a concrete example, consider
\[
\Axiom{\F{Expr}\hab\CTR}\quad
\F{Expr} = \ctr{\ENUM{\C{neutral},\C{combine}}}
  {\left\{\begin{array}[c]{@{}r@{\;\mapsto\;}l}
    \C{neutral} & \ENUM{} \\
    \C{combine} & \ENUM{\C{left},\C{right}}
          \end{array}\right.}
\]
which induces $\FM{\F{Expr}}{\V{X}}$, the type of terms in the \remph{syntax} of monoids
with variables drawn from $\V{X}$.

Notoriously, it also takes function extensionality to even prove that the neutral element
\[
\nd{\pr{\C{neutral}}{\F{absurd}}}
\qquad\mbox{or}\qquad
\nd{\pr{\C{neutral}}{\F{const}\:\V{t}}}
\]
is unique, because, on the face of it, intensionally it is not. There are lots of ways to
construct functions from the empty type to expressions --- e.g., by choosing any $t$, above.
And if we were so unwary as to identify all of these functions \jgm{ally}, we should
collapse the \jgm{al} equality in the presence of false hypothesis and destroy the computability
of normal forms.

We pay quite a price for the functional encoding of first-order things! Let us not!

Our plan is to replace $\SET$ in the definition of $\CTR$ by our universe $\UF$ of finite types,
so that the $\fto$ in the definition of $\CX{\cdot}{\cdot}$ can be replaced by the \remph{tabulation} of the function space from a finite position type.

\section{Finite Types and our First Attempt at Tabulated Functions}

We intend to construct a universe $\UF$ of finite types with an explicit interpretation%
\footnote{hence a \emph{Tarski}-style universe} into $\SET$.
\[
\Axiom{\UF\hab\SET}\qquad
\Rule{\V{N}\hab\UF}
     {\ElF{\V{N}}\hab\SET}
\]
A key objective, however, is also to obtain a tabulated dependent function type, whose domain is finite but whose codomain need not be
\[
\Rule{\V{N}\hab\UF \hg \V{T}\hab\ElF{\V{N}}\fto\SET}
     {\VERTPI{\V{N}}{\V{T}}\hab \SET}
\qquad
\Rule{\V{f}\hab \PI{x}{\V{N}}{\V{T}\:\V{x}}}
     {\vertflam{\V{f}}\hab \VERTPI{\V{N}}{\V{T}}}
\qquad
\Rule{\V{ts}\hab \VERTPI{\V{N}}{\V{T}} \hg \V{i}\hab\ElF{\V{N}}}
     {\vertfapp{\V{ts}}{\V{i}}\hab \V{T}\:\V{i}}
\]
\noindent for which a $\beta$-rule is provable, namely
\[\vertfapp{\vertflam{\V{f}}}{\V{i}}\peq\V{f}\:\V{i}\]
\noindent where attempting to have this rule hold \jgm{ally} will get us into the same exact trouble as before.

\paragraph{What do we want in $\UF$?} We shall certainly need finite enumerations, such as the
$\ENUM{\C{neutral},\C{combine}}$ and $\ENUM{\C{left},\C{right}}$ encountered above. It would helpful
if these were not both the \emph{number} $2$. Moreover, to facilitate constructions such as the derivative
of containers, we should support dependent pairing: the shape of the derivative of $\ctr{\V{S}}{\V{P}}$
is the dependent pair type $ \SG{s}{\V{S}}{\V{P}\:\V{s}}$. For our $\F{Expr}$ example, these shapes are
$(\pr{\C{combine}}{\C{left}})$ and $(\pr{\C{combine}}{\C{right}})$ which, by an outrageous coincidence,
also add up to $2$. It should surely be possible to extract the original shape components without resorting
to arithmetic. The ergonomic effectiveness of the presentation is more than the numerical adequacy of the 
representation.

We can satisfy our latter concern in short order by closing $\UF$ under dependent pairing, decoded as
the dependent pairing of the meta-language.
\[
\Rule{\V{S}\hab\UF \hg \V{T}\hab\ElF{\V{S}}\fto\UF}
     {\FSG{\V{S}}{\V{T}}\hab \UF}
\qquad
\ElF{\FSG{\V{S}}{\V{T}}} = \SG{s}{\V{S}}{\V{T}\:\V{s}}
\]
We promise to return to the representation of finite tagged enumerations in due course: indeed $\UF$ will need at least one base case. However, in choosing the narrative structure for this story, we prefer to wait until we have explored the consequences of our choice to close $\UF$ under dependent pairing. We were taught a lesson!

Our presentation of tabulated functions for dependent pairs is simply currying
\[
\VERTPI{(\FSG{\V{R}}{\V{S}})}{\V{T}}\; =
\; \VERTPI{\V{R}}{\LAM{r}{\VERTPI{\V{S}\:\V{r}}{\LAM{s}{\V{T}\:(\pr{\V{r}}{\V{s}})}}}}
\]
where our convention is that $\LAM{\cdot}{}$ scopes rightward as far as possible and
that $\vertfapp{\cdot}{\cdot}$ associates to the left%
\footnote{``rather as we all did in the Sixties'' -- Roger Hindley}
\[ \vertflam{\V{f}}\; =\; \vertflam{\LAM{r}{\vertflam{\LAM{s}{\V{f}\:(\pr{\V{r}}{\V{s}})}}}}
\qquad
\vertfapp{\V{f}}{(\pr{\V{r}}{\V{s}})}\; =\; \vertfapp{\vertfapp{\V{f}}{\V{r}}}{\V{s}}
\]
Our careful reader will regard these equations as suspicious, with no clue in the code as to the
type-level recursion which is unfolding. Quite right, too! 

Let us fill in the hidden $\UF$-codes that clarify the type-level recursion in the above.
\[
\begin{array}{rcl}
\vertflam{\impl{\FSG{\V{R}}{\V{S}}}\:\V{f}}& = &
  \vertflam{\impl{\V{R}}\:\LAM{r}{\vertflam{\impl{\V{S}\:\V{r}}\:\LAM{s}{\V{f}\:(\pr{\V{r}}{\V{s}})}}}} \\
\vertfapp{\V{f}}{\impl{\FSG{\V{R}}{\V{S}}}\:(\pr{\V{r}}{\V{s}})} & = & 
  \vertfapp{\vertfapp{\V{f}}{\impl{\V{R}}\:\V{r}}}{\impl{\V{S}\: \V{r}}\:\V{s}}
\end{array}
\]

% should we have a (code for) internal Pi types?
What if the codomain of such a function is also finite? Such functions are themselves finitely enumerable. Thus we might wish to have them in $\UF$ somehow, as well. First, we can replay the previous construction of tabulated function types entirely within $\UF$, obtaining

\[
\Rule{\V{N}\hab\UF \hg \V{T}\hab\ElF{\V{N}}\fto\UF}
     {\VERTPIF{\V{N}}{\V{T}}\hab \UF}
\]
Note that we have move the $\UF$ subscript to the codomain end of the arrow. Abstraction and application are exactly as before. Of course, we could have written $\VERTPI{\V{N}}{\LAM{n}{\ElF{\V{T}\:\V{n}}}}$ if we were content with a $\SET$, but we are not. The point is to get function spaces \emph{internal} to $\UF$. This rather raises the question of whether we should go further and add codes for function spaces as a constructor for $\UF$. The good news is that we may readily take 
\[ \ElF{\ROUGEPIF{\V{N}}{\V{T}}}\; = \; \VERTPI{\V{N}}{\LAM{n}{\ElF{\V{T}\:\V{n}}}} \]
using our tabulated function space to avoid the function space of our metalanguage. The bad news is that we now have to extend this tabulated function space construction to the case when the domain is itself a function type:
\[
\VERTPI{(\ROUGEPIF{\V{R}}{\V{S}})}{\V{T}}\; = \; ?
\]
We can see no direct way to compute this higher-order function type more neatly than to tabulate its domain and then curry, so adopting the constructor for function spaces does not save us from computing them. We face less work without it. We may not have $\UF$-codes for the function spaces but we can make do with the $\UF$-codes for their tabulations.
% sadly 'tetration' is not categorical...

The price we pay is that $\ElF{\VERTPIF{\V{N}}{\V{T}}}$ is \emph{not} the same type as $\VERTPI{\V{N}}{\LAM{n}{\ElF{\V{T}\:\V{n}}}}$ when $\V{N}$ and $\V{T}$ are open terms. Equating these types makes sense \remph{algebraically} but does not follow from the definitions. It has the potential, however, to be adopted as a `$\nu$-rule' in the sense of \citet{DBLP:conf/icfp/AllaisMB13}, as it is sufficient to add it only for open terms which have stopped computing by pattern matching.

\paragraph{Finitary Containers and their Free Monads}
Now that we have tabulated functions, let us adapt containers and see what happens to the free monad construction. ($\KON{\blank}$ makes constant functions.)
\[
\Rule{\V{S}\hab\UF\hg \V{P}\hab\ElF{\V{S}}\fto\UF}
     {\ctr{\V{S}}{\V{P}}\hab\FTR}
\qquad
\Rule{\V{C}\hab\FTR\hg \V{X}\hab\SET}
     {\FX{\V{C}}{\V{X}}\hab\SET}
\quad
\FX{\ctr{\V{S}}{\V{P}}}{\V{X}} = \SG{s}{\ElF{\V{S}}}{\VERTPI{\V{P}\:\V{s}}{\KON{\V{X}}}}
\]
Note that we are quite happy to see \remph{actual} functions in the \remph{descriptions} of containers, but we insist on \remph{tabulated} functions in their extensions.
\[
\Rule{\V{C} \hab \FTR \hg  \V{X} \hab \SET}
     {\FM{\V{C}}{\V{X}} \hab \SET}
\qquad
\Rule{\V{x}\hab \V{X}}
     {\va{\V{x}}\hab \FM{\V{C}}{\V{X}}}
\qquad
\Rule{\V{cts} \hab \FX{\V{C}}{\pCX{\FM{\V{C}}{\V{X}}}}}
     {\nd{\V{cts}}\hab \FM{\V{C}}{\V{X}}}
\]
Now, our term-like data are genuinely first-order things, with just one implementation each. But what happens to join?
\[ \Rule{\V{ts} \hab \FM{\V{C}}{(\FM{\V{C}}{\V{X}})}}
        {\F{join}\: \V{ts} \hab \FM{\V{C}}{\V{X}}}
\qquad\qquad
\begin{array}[c]{lcl}
\F{join}\: \va{\V{x}} & = & \V{x}\\
\F{join}\: \nd{\pr{\V{s}}{\V{k}}} & = & \nd{\pr{\V{s}}{\vertflam{\LAM{p}{\F{join}\: (\vertfapp{\V{k}}{\V{p}})}}}} \\
\end{array}
\]
This definition looks good on paper but arouses the suspicions of most syntactically driven termination checkers because $\vertfapp{\V{k}}{\V{p}}$ is not conspicuously smaller than $\nd{\pr{\V{s}}{\V{k}}}$, even though it is in every case projected from $\V{k}$. We could, perhaps, work around this problem by laboriously inlining the tabulation and projection operations, but why should we have to? The power of the container abstraction is that we explain how to implement operations like $\F{join}$ \remph{uniformly}. Can we construct a recursion \remph{explanation} which justifies the uniform approach?


\section{The View from Above, Obscured}

A lesson well worth learning from \citet{DBLP:journals/jar/McKinnaP99} is that data can admit newer and more convenient induction principles than those with which they were conceived. We may reorganise the recursive structure of data in much the way that \citet{DBLP:conf/popl/Wadler87} rearrange case analysis with `views'. Let us say what it is for a term in a finitary free monad to be `recursive' in the sense that its notion of substructure projection is well founded.
\[
\Rule{\V{t} \hab \FM{\V{C}}{\V{X}}}
     {\REC{\V{t}} \hab \SET}
     \qquad
\Rule{\V{x}\hab \V{X}}
     {\va{\V{x}}\hab \REC{\va{\V{x}}}}
\qquad
\Rule{\V{s}\hab\ElF{\CSh{\V{C}}} \quad
      \V{f}\hab\PI{p}{\CPo{\V{C}}{\V{s}}}{\REC{(\vertfapp{\V{k}}{\V{p}})}}}
     {\nd{\pr{\V{s}}{\V{f}}} \hab \REC{\nd{\pr{\V{s}}{\V{k}}}}}
\]
Note that our functional presentation in the step case of recursivity makes $\V{f}\:\V{p}$ clearly a structural subproof of $\nd{\pr{\V{s}}{\V{f}}}$. Correspondingly, we achieve projection-based recursion on a term by ordinary structural recursion on its recursivity proof. The challenge now is to implement the `covering function', in the sense of \citet{DBLP:journals/jfp/McBrideM04}, which shows that all terms are recursive. And that's when our troubles \remph{really} begin.

Our task is to implement, using only structural recursion,
\[ 
\Rule{\V{t} \hab \FM{\V{C}}{\V{X}}}
     {\rec{\V{t}} \hab \REC{\V{t}}}
\]

Our opening moves are forced, the base case is straightforward and in the step case we must clearly copy the shape across and receive a position by return:
\[ 
\begin{array}[c]{lcl}
\rec{\va{\V{x}}} & = & \va{\V{x}}\\
\rec{\nd{\pr{\V{s}}{\V{k}}}} & = & \nd{\pr{\V{s}}{\LAM{p}{\HOLE}}} \\
\end{array}
\]
At this point, we see that we need
\[
\V{k} \hab \VERTPI{\CPo{\V{C}}{\V{s}}}{\KON{(\FM{\V{C}}{\V{X}})}} \quad 
\vdash \quad
\HOLE \hab \REC{(\vertfapp{\V{k}}{\V{p}})}
\]
but we are not allowed simply to \remph{apply} $\V{k}$. Rather we must dismantle it carefully according to the structure of $\CPo{\V{C}}{\V{s}}$. So we shall certainly need a helper function which abstracts that position type in order to inspect it.

\[
\Rule{\V{P}\hab\UF\quad \V{k} \hab \VERTPI{\V{P}}{\KON{(\FM{\V{C}}{\V{X}})}}
      \quad \V{p}\hab \ElF{\V{P}}}
     {\F{help}\:\V{P}\:\V{k}\:\V{p} \hab \REC{(\vertfapp{\V{k}}{\V{p}})}}
\]
The difficult case is when $\V{P}$ is a dependent pair type.
\[
\F{help}\:(\FSG{\V{R}}{\V{S}}) \:\V{k}\:(\pr{\V{r}}{\V{s}}) = \HOLE
\quad\text{where}\quad
\V{k} \hab 
\VERTPI{\V{R}}{\LAM{r}{\VERTPI{\V{S}\:\V{r}}{\KON{(\FM{\V{C}}{\V{X}})}}}}
\]
We certainly would like to call $\F{help}\:\V{S}\:\HOLE\:\V{s}$ but the return type of $\V{k}$ is not helpful for $\HOLE$ because the terms we are trying to unpack are now two layers down instead of one. This problem will only get worse if $\V{R}$ turns out also to be a dependent pair type, pushing the term even further down. We have an entire class of problems to solve by mutual recursion and so far we have seen only three instances of it.

\paragraph{Dependent types to the rescue!} The dependently typed way to solve a class of problems is to \remph{reflect} it. We started with the problem to show the recursivity of a single term but our problems kept growing by branching over position sets.

We can instantiate a general pattern for solving such problems. Step one, introduce a \remph{data} type $\D{Problem}\hab\SET$ of problem \remph{descriptions}. Step two, recursively define $\F{Input}\hab \D{Problem}\fto\SET$ which specifies the input data appropriate to the problem in question. Step three, recursively define $\F{Output}\hab \PI{P}{\D{Problem}}{\F{Input}\: \V{P}\fto\SET}$. Step four, solve all the problems by defining $\F{solve}\hab\PI{P}{\D{Problem}}{\PI{i}{\D{Input}\:P}{\F{Output}\:P\:i}}$. Crucially $\V{P}$ can swell as well as shrink while $\F{solve}$ proceeds provided $\V{i}$ gets smaller.

\newcommand{\basecase}{\C{\varepsilon}}
\newcommand{\branch}[2]{#1\mathop{\C{\Yleft}}#2}
In our case, we should start with a base case which shows where we come in
\[\Axiom{\basecase \hab \D{Problem}}
\qquad
\begin{array}[c]{l}
\F{Input}\: \basecase \;=\; \FM{\V{C}}{\V{X}}
\\
\F{Output}\: \basecase\: \V{t} \;=\; \REC{\V{t}}
\end{array}
\]
so that we can take $\rec{\V{t}} = \F{solve}\:\basecase\: t$. Then our step case, $\branch{\blank}{\blank}$, should allow us to branch problems over a finite set.
\[
\Rule{\V{S}\hab\UF \quad \V{P}\hab \ElF{S}\fto\D{Problem}}
     {\branch{\V{S}}{\V{P}} \hab \D{Problem}}
\quad
\begin{array}[c]{l}
\F{Input}\: (\branch{\V{S}}{\V{P}}) \;=\; \VERTPI{\V{S}}{\LAM{s}{\F{Input}\:(\V{P}\:\V{s})}}
\\
\F{Output}\: (\branch{\V{S}}{\V{P}})\: \V{k} \;=\; 
    \PI{s}{\ElF{\V{S}}}{\F{Output}\:(\V{P}\:\V{s})\:(\vertfapp{\V{k}}{\V{s}})}
\end{array}
\]

With this extra flex in problem description, let us now try to $\F{solve}$ our whole family of problems.
\[\begin{array}{lcl}
\F{solve}\:\basecase\:\va{\V{x}} & = & \va{\V{x}}\\
\F{solve}\:\basecase\:\nd{\pr{\V{s}}{\V{k}}} & = &
\nd{\pr{\V{s}}{\F{solve}\:(\branch{\CPo{\V{C}}{\V{s}}}{\KON{\basecase}})\:\V{k}}} \\
\F{solve}\:(\branch{(\FSG{\V{R}}{\V{S}})}{\V{P}})\:\V{k} & = &
\F{solve}\:
(\branch{\V{R}}{\LAM{r}{
  \branch{\V{S}\:\V{r}}{\LAM{s}{\V{P}\:(\pr{\V{r}}{\V{s}})}}}})
\:\V{k}\\
\vdots\\
\end{array}\]
Of course, we shall need to spell out our story about enumerations to see how the mutual recursion returns from $\branch{\blank}{\blank}$ problems to $\basecase$ problems, but the good news is that currying out the problem descriptions syncs neatly with the currying out of the type of $\V{k}$ and the story so far typechecks. The bad news, however, is that the problem acquires a larger sequence of smaller branchings while $\V{k}$ stays put, so syntactic guard checkers are fooled into rejecting this entirely sensible definition.


\section{Bleu et Vert: Tabulated Functions Mille Feuilles}

When all you have is a hammer, everything looks like a nail. Conspicuous constraction comes from cancelling construction. Our currying by computing recursively over $\UF$ has left the after currying $\V{k}$ look no simpler than the $\V{k}$ before currying. Excessive currying conceals the subtle layering of problem simplification that is our intent. A more sophisticated cuisine is required.

We can arrest computation and make simplification explicit by discarding our green-typeset tabulated function space $\VERTPI{\V{N}}{\V{P}}$ and reintroducing the notion \remph{inductively}, with a canonical (therefore, typeset in blue) type constructor and an explicit value constructor.
\[
\Rule{\V{N}\hab\UF \hg \V{T}\hab\ElF{\V{N}}\fto\SET}
     {\BLEUPI{\V{N}}{\V{T}}\hab \SET}
\qquad
\Rule{\V{k}\hab\VERTPIB{\V{N}}{\V{T}}}
     {\nd{\V{k}}\hab\BLEUPI{\V{N}}{\V{T}}}
\]
But what is being wrapped by $\nd{\blank}$? Exactly a \remph{small-step} version of our computed tabulation!
\[
\Rule{\V{N}\hab\UF \hg \V{T}\hab\ElF{\V{N}}\fto\SET}
     {\VERTPIB{\V{N}}{\V{T}}\hab \SET}
\qquad     
\VERTPIB{(\FSG{\V{R}}{\V{S}})}{\V{T}} \;=\;
\BLEUPI{\V{R}}{\LAM{r}{\BLEUPI{\V{S}\:\V{r}}{\LAM{s}{\V{T}\:(\pr{\V{r}}{\V{s}})}}}}
\]
This function does not compute \remph{recursively}. Rather, after one step of case analysis, it reinvokes the \remph{inductive} tabulation (in strictly positive positions, of course!). This is a reusable culinary technique, mutually defining inductive and small-step computational versions of the same notion, to brake the unfolding of computation by requiring the unwrapping of data. We call this delicately layered recipe `bleu et vert'.

Now, we must repair \remph{everything}, changing the types of abstraction and application to their `bleu' counterparts. Note that both operations remain \remph{defined} themselves, so their names stay typeset in green.
\[
\Rule{\V{f}\hab \PI{x}{\V{N}}{\V{T}\:\V{x}}}
     {\bleuflam{\V{f}}\hab \BLEUPI{\V{N}}{\V{T}}}
\qquad
\Rule{\V{ts}\hab \BLEUPI{\V{N}}{\V{T}} \hg \V{i}\hab\ElF{\V{N}}}
     {\bleufapp{\V{ts}}{\V{i}}\hab \V{T}\:\V{i}}
\]
The definitions must still be constructed with some explicit awareness of type descriptions, but now we must also be careful to wrap and unwrap in the right places.     
\[
\begin{array}{rcl}
\bleuflam{\impl{\FSG{\V{R}}{\V{S}}}\:\V{f}}& = &
  \nd{\bleuflam{\LAM{r}{\bleuflam{\LAM{s}{\V{f}\:(\pr{\V{r}}{\V{s}})}}}}} \\
\bleufapp{\nd{\V{f}}}{\impl{\FSG{\V{R}}{\V{S}}}\:(\pr{\V{r}}{\V{s}})} & = & 
  \bleufapp{\bleufapp{\V{f}}{\V{r}}}{\V{s}}
\end{array}
\]
Specifically, we must match explicitly on type descriptions to determine \remph{what} is to be wrapped by abstraction or unwrapped by application. At use sites, including on the right hand sides of the definitions, the use of a canonical type constructor for tabulated function space allows the hidden type information to be reconstructed unambiguously. This brevity is an unexpected bonus from the layers of bleu separating the vert.

\paragraph{The `Bleu et Vert' View from Above}
We should now check that our extra layering makes our recursion structural. Our `$\D{Problem}$, $\F{Input}$, $\F{Output}$, $\F{solve}$' approach remains the same, apart from the step case of $\F{Input}$, which crucially now becomes
\[
\F{Input}\: (\branch{\V{S}}{\V{P}}) \;=\; \BLEUPI{\V{S}}{\LAM{s}{\F{Input}\:(\V{P}\:\V{s})}}
\]
giving us the layers we need to be seen to strip away. We obtain
\[
\F{solve}\:(\branch{(\FSG{\V{R}}{\V{S}})}{\V{P}})\:\nd{\V{k}}  \;=\;
\F{solve}\:
(\branch{\V{R}}{\LAM{r}{
  \branch{\V{S}\:\V{r}}{\LAM{s}{\V{P}\:(\pr{\V{r}}{\V{s}})}}}})
\:\V{k}
\]
which is clearly a simplification.

% that typechecks but doesn't termination check. SIGH.

%HERE
% We should of course now give a treatment of finite enumerations, but first let us see how our free monads with tabulated functions plan is going. If in join we replace (k p) with (k $F p) we are no longer structurally recursive! So internalisation goes blam. Moreover if we try to inline \lambda F and $F, that forces us to generalise the problem to cope with currying. We very quickly lose sight of why the recursion is (actually) structural. Worse, we are doing a complicated maneuver one recursive function over the data. That insures our approach will not scale, as we do not want to be forced to do this every time. We want to introduction a recursion view and we need some infrastructure to deal with that (bleu/vert). Once we've figured our way out of that mess, enumerations will come out `right' on their own.

% Worth commenting? The bleu/vert idea shows up in many more places, including in symbolic computation. And gluing.


\hrulefill
\begin{itemize}
\item leverage that finiteness: we want a universe that guarantees finiteness and decidability of 
some questions of an exploratory nature, which are decidable when exploring is bounded.
\item also motivates tabulated functions
\end{itemize}

Points to make:
\begin{enumerate}
\item finite types as coordinate system for positions in data types
\item when finite types are domains of functions, those functions have first-order tabulations
\item to be clear: Universes, DataEqDec, Equality not used in the story directly
% \item it's useful to have at least 2 universes of proofs 
% \item being polymorphic over all of that... doesn't seem to work very well; which unfortunately
%  leads to some duplication (see FiniteEq amongst others)
\item initerested in (structured!) finite types as presentations of evidenced algorithmic decision making
\item bleu / vert; internal / external
% motivation 1: bleu ~ invertibility of typechecking. leads to winning with the termination checker
% purely vert, you lose
% post-facto: intension/extension, representation/computation pairings
\item do coerce and coherence simultaneously (not crucial, but ergonomic)
\item want to talk about 2 positions being equal (and the evidence of that) without leaving UF,
  i.e. why it is important UD is interpreted by UF
\item "classical" logic becomes constructive for UD+UF. That's worth a lot.
\end{enumerate}

Minor things to remember:
\begin{enumerate}
\item adventures in beating the termination checker
\item Pi types in finite types are no fun
\item Agda-ism: (proof) irrelevance for things you know are irrelevant seems to work;
 proof irrelevance as a \emph{desiderata} is hellish
\item Strings are so much nicer than numbers... and still finite. Humans like them, especially for tags.
\item where \emph{do} we use thinnings?
\item heterogeneous equality, and when off-diagonal should be `0 vs `1
\item J contains choices
\end{enumerate}

Story beats:
\begin{enumerate}
\item Introduction
\item Free Monads of Containers

\item Infrastructure to support the above
\begin{itemize}
\item universe F of finite types; key bit: finite enumerations E
\item we want some structure on top of that
\item do we want Pi? we certainly don't want Agda's functions as what we interpret anything in E
\item internal tabulated functions: lo and behold, they are vert!
\item next: external tabulated functions or nice E? choose nice E.
\end{itemize}

\item Building Free Monads
\begin{itemize}
\item interested in tabulated functions where codomain is not finite
\item want to build recursive view that gives us successive layers (ordinal interface) of those
structures
\item the view should cover. the termination checker does not believe us. It's not wrong, as
it can't get its hands on the evidence that it actually is.
%tabulated functions: there are ways to do it that will lead to big fights with the termination
% checker. See comments above. open point: which order? % TODO
\item reminder: you might think you want recursivity evidence to be proof irrelevant... but no, you don't.
yes, they are mere propositions, but you need to be strict lest you enable the liars to win
\item with the recursive view actually in hand, we can get there, and prove the Monad laws honestly
\item use inductive equality here, because we're exactly dealing with first-order data that is all Set-like
\item since substitution are (Agda) functions, the statement of the laws need to be stated 'extensionally'
\end{itemize}

\item Searching Free Monads
\begin{itemize}
\item decisions, De Morgan, dneg
\item decidable equality, very classical, have LEM, (X -> False) -> False goes away
% we could get away with homogeneous, with proper book-keeping. but that's a story for another day.
\item important point: no quadratic search. Testing equality of shapes unifies position types (J!)
\item one-hole context / pointing. every value has a finite coordinate system (either way)
\item decisions about values lift because of the precision of the coordinate system
\end{itemize}
\end{enumerate}

TODO: inverse image (because we can search!)

Potential end: Many are obsessed with (very) large universes. We ought to show some love to the small ones.
UD+UF just the start.

\bibliography{searchme}

\end{document}

